\section{Related Work}

\subsection{Agile Traceability and Requirement Traceability}

Requirement traceability connects requirements with related artifacts across the
software lifecycle. Prior work shows that natural language processing, information
retrieval, and learning-based techniques can help recover or maintain trace links
between textual artifacts \cite{nlpTraceability,crossLevelTraceability}. Other work
uses bag-of-words, TF-IDF, word embeddings, BERT-based models, and retrieval-augmented
generation to improve semantic matching and trace link recovery
\cite{crossLevelTraceability,traceBert,ragTraceability}.

These studies are useful for long-term automation, but they often focus on recovering
links after artifacts already exist. DevTrack AI takes a more operational direction
for student teams. It keeps traceability visible during project execution through a
rule-based RTM and human-confirmed artifact links. AI-based trace suggestions can be
added later, but official traceability evidence should remain explainable and
reviewable in the academic context.

\subsection{Requirement, Test, Defect, and Evidence Management}

Traceability is useful only when connected artifacts are meaningful and current. A
requirement should be related not only to implementation tasks, but also to acceptance
criteria, test cases, defects, and proof of completion. Requirement-driven tracking
therefore needs links between planning artifacts and verification artifacts so that
project status can be explained from evidence \cite{nlpTraceability,crossLevelTraceability}.

In student projects, this relationship is often weak. Teams may write test cases
late, record defects informally, or upload screenshots without review. DevTrack AI
addresses this gap through an Evidence Vault and a review-oriented task completion
flow. Evidence can include screenshots, API links, test results, GitHub issues,
commits, pull requests, and CI/check metadata. A task may request completion review,
but the leader still approves or rejects it based on the available evidence. This
keeps traceability useful for monitoring and collaboration, not only final
documentation \cite{ragTraceability}.

\subsection{Automated Reporting and Software Artifact Summarization}

Automated reporting is related to research on release note generation, pull request
description generation, and commit message generation. These studies commonly treat
software development data as source material for summarization. For example,
extractive summarization can select important commits or pull requests for release
notes \cite{releaseNotes}. Pull request description generation studies show how
commit messages and source-code context can be transformed into natural language
descriptions \cite{prDescriptions}. Commit message generation research similarly explores
how source-code changes can be summarized into readable descriptions
\cite{commitMessages}.

The lesson for DevTrack AI is that weekly reports should not be generated from an
empty prompt. The system first collects structured facts from tasks, sprint status,
RTM snapshots, defects, evidence review status, and GitHub activity. AI can then
rewrite those facts into a readable draft. This keeps the report grounded in project
data and allows the leader to review it before it is saved or submitted.

\subsection{Mining Software Repositories and AI for Software Engineering}

Mining Software Repositories (MSR) studies how data from version control systems,
issue trackers, pull requests, and other development platforms can be used to
understand software development activities \cite{teachingMsr}. This is relevant to
DevTrack AI because code contribution is an important part of student project
assessment, but commit count alone is not a fair indicator of contribution. DevTrack
AI therefore combines linked commits, pull requests, CI/check results, task
completion, accepted evidence, test activity, defect fixes, and review decisions.

Large language models (LLMs) have also been studied for software engineering tasks
such as code summarization, code understanding, issue analysis, and development
support \cite{llmSurvey}. In project-based
courses, LLM-based contribution summarization is also relevant because it can convert
repository activity into a more understandable description of individual work
\cite{llmContribution}. However, LLM output must be controlled because it can
provide plausible but unverifiable explanations. In DevTrack AI, AI support is
therefore used for explanation and drafting, not for automatic grading or task
approval. This human-in-the-loop principle is important for academic transparency.
